\section*{الفصل \ref{ch:b3:conservation-biology} --- بيولوجيا الصون}

\begin{solution}{exo:b3:conservation-biology:1}
\omterm{def:b3:conservation-biology:small}{العشوائية الديموغرافية}: المصادفة في مواليد بضعة أفراد
ووفياتهم وأجناسهم --- الكاكابو عند واحد وخمسين، بذكور أكثر من
الإناث. \omterm{def:b3:conservation-biology:small}{والعشوائية البيئية}: سنة سيئة تضرب كل فرد
دفعةً واحدة --- جفاف، أو مرض، أو شتاء قاس يصيب
جمهرة من بضع مئات. \omterm{def:b3:conservation-biology:small}{وأثر ألي}: هبوط النمو للفرد
عند الكثافة المنخفضة --- الحمام الزاجر، الذي لم تقدر مستعمراته على التكاثر
متى رقّت، والذي لم تعد مفترساته الباقية
تُشبَع.
\end{solution}

\begin{solution}{exo:b3:conservation-biology:2}
حجم جمهرة مثالية --- ثابتة، متزاوجة عشوائيًا، متساوية
الجنسين، بأحجام أسر بواسونية --- تفقد \omterm{thm:b3:conservation-biology:ne}{التغايرية الزيجية}
بالمعدل نفسه الذي تفقدها به الحقيقية. وأما السمات: فعدد غير متساو من الذكور
المتكاثرين والإناث المتكاثرات؛ وتقلبات في الحجم، ترجّح الأجيال الأصغر؛
وتباين في حجم الأسرة فوق البواسوني، إذ ينتج بضعة أفراد
معظم الصغار (وكذلك التزاوج غير العشوائي وتداخل
الأجيال).
\end{solution}

\begin{solution}{exo:b3:conservation-biology:3}
$S = cA^{z}$. والنسبة المحفوظة $= (A'/A)^{z}$: أي $0.5^{0.25} = 0.84$؛
و$0.1^{0.25} = 0.56$؛ و$0.01^{0.25} = 0.32$ --- أي إن جزءًا من مئة من المساحة
يبقي مع ذلك ثلث الأنواع، في النهاية.
\end{solution}

\begin{solution}{exo:b3:conservation-biology:4}
$N_{e} \ge 50$ يبقي ارتفاع التزاوج الداخلي دون $1/100$ في
الجيل، وهو المعدل الذي يجد مربّو الحيوان التدهور عنده
محتمَلًا: فهو يحمي من فقد قصير المدى للملاءمة. و$N_{e} \ge
500$ يوازن فقد التباين بالانجراف بتزويد الطفر
له، فتحفظ الجمهرة التباين الكمي الذي تحتاج إليه
للتأقلم: فهو يحمي المدى الطويل. والجمهرات الحقيقية تحتاج إلى خمسة إلى عشرة
أضعاف هذين العددين في التعداد.
\end{solution}

\begin{solution}{exo:b3:conservation-biology:5}
$N_{e} = 4\times 8\times 120/128 = 30$؛ وتهبط \omterm{thm:b3:conservation-biology:ne}{التغايرية الزيجية} بمقدار $1/60 =
\qty{1.7}{\%}$ في الجيل؛ ويُفقد ربعها حين يكون $(1 - 1/60)^{t} =
0.75$، أي $t = \ln 0.75/\ln(59/60) = 17$ جيلًا.
\end{solution}

\begin{solution}{exo:b3:conservation-biology:6}
$1/N_{e} = \tfrac{1}{5}(4/1000 + 1/20) = 0.0108$، أي $N_{e} = 93$: فجيل
واحد عند عشرين جعل خمسة أجيال تسلك سلوك ثلاثة وتسعين.
\omterm{thm:b3:conservation-biology:ne}{والتغايرية الزيجية} المحفوظة $(1 - 1/186)^{5} = 0.973$، في مقابل $(1 -
1/2000)^{5} = 0.9975$ لألف ثابتة.
\end{solution}

\begin{solution}{exo:b3:conservation-biology:7}
$F = 1 - (1 - 1/20)^{5} = 0.23$؛ وتهبط الملاءمة $5\times 2.3 = \qty{11}{\%}$.
وكل خلف لأنثى من تكساس وذكر من فلوريدا عنده $F = 0$؛ فإذا كانت
الوافدات الثماني خُمسَي الإناث المتكاثرة، كان خُمسا
الجيل التالي من تزاوج بعيد وهبط متوسط $F$ إلى نحو
$0.6\times 0.23 = 0.14$ في جيل واحد --- وهو ما أظهره
التعافي.
\end{solution}

\begin{solution}{exo:b3:conservation-biology:8}
$r = \ln(100/31)/8 = \qty{0.15}{yr^{-1}}$. ولو استمر إلى 2020:
$100\,\mathrm{e}^{0.15\times 17} \approx 1200$. والذي حواه المتنزه نحو
مئة: فقد امتلأت الأقاليم، وتراجعت الأيائل، وقتلت القطعان أفراد
بعضها بعضًا، وانتشر الجرب والسُّلاب --- أي إن الاعتماد على الكثافة بدأ
متى شُغل المدى الفارغ.
\end{solution}

\begin{solution}{exo:b3:conservation-biology:9}
(أ) تراجع قدره \qty{60}{\%} في عشر سنين يتجاوز \qty{50}{\%}: مهدَّد بالانقراض.
(ب) $180 < 250$ فردًا ناضجًا: مهدَّد بالانقراض تهديدًا حرجًا. (ج) مدى
\qty{3000}{km^2}، دون \qty{5000}{}: مهدَّد بالانقراض. (د) \num{30000}
فرد وتراجع قدره \qty{10}{\%}: لا عتبة منهما بُلغت --- فليس
مهدَّدًا (وأقصاه قريب من التهديد، يُراقَب).
\end{solution}

\begin{solution}{exo:b3:conservation-biology:10}
تخفق الأزواج $N$ كلها باحتمال $(1/2)^{N}$: $0.25$، و$0.031$،
و$0.001$، و$10^{-6}$. فالخطر الديموغرافي يهبط أسّيًا مع $N$ وهو
مهمَل بعد بضع عشرات من الأفراد؛ وفوق ذلك يكون الخطر
الباقي بيئيًا، وهو لا يهبط بتلك السرعة، لأنه
يضرب كل الأفراد معًا.
\end{solution}

\begin{solution}{exo:b3:conservation-biology:11}
شظية واحدة من عُشر: $0.1^{0.25} = 0.56$ من الأنواع الأصلية.
وعشر شظايا تحوي بين $0.56$ (إن حوت كلها الأنواع نفسها) و،
مبدئيًا، أكثر من الكل (إن حوت كل واحدة أنواعًا مختلفة ---
وهو مستحيل فوق المجمع الإقليمي). وأين تقع يتوقف على
التبدل بين الشظايا: فاختلاف الموائل بين الشظايا
يدفع المجموع إلى الأعلى؛ والأنواع التي تحتاج إلى مساحات كبيرة أو موئل داخلي، أو التي
تُفقد من كل شظية بآثار الحافة نفسها، تدفعه إلى
الأسفل؛ ودين الانقراض في كل شظية يُسدَّد على حدة، فيكون
مجموع المدى الطويل عادةً دون مجموع الكل.
\end{solution}

\begin{solution}{exo:b3:conservation-biology:12}
يزيل الانجراف \omterm{thm:b3:conservation-biology:ne}{التغايرية الزيجية} بمعدل $1/2N_{e}$ في الجيل؛
ويستبدل المهاجرون نسبة $m$ من مجمع المورّثات بأليلات مسحوبة
من مكان آخر. وعند $mN_{e} \approx 1$ يكون $m \approx 1/N_{e}$، أي ضعف
معدل الفقد، فيُجدَّد التباين أسرع مما يضمحل وتبقى
الجمهرة قريبة من المصدر في تواتر الأليلات ($F_{ST} = 1/(1
+ 4N_{e}m) \approx 0.2$). والكلفة: تُخفَّف الأليلات المحبَّذة موضعيًا
بمقدار $m$ في الجيل، فلا يبقى التأقلم الموضعي إلا لصفات
يتجاوز معامل انتخابها نحو $1/N_{e}$؛ وأما الفوارق الموضعية
الضعيفة الانتخاب فتُغمر.
\end{solution}

\begin{solution}{pb:b3:conservation-biology:1}
\textbf{1.} $N_{e} = 4\times 20\times 40/60 = 53$، في مقابل تعداد قدره 60.
\textbf{2.} $1/N_{e} = \tfrac{1}{5}(1/500 + 1/100 + 1/60 + 1/51 + 1/60) =
0.0130$، أي $N_{e} = 77$؛ ويهيمن جيلا عنق الزجاجة، 51 و60،
ولا تكاد الخمسمئة تُحسب.
\textbf{3.} $F = 1 - (1 - 1/154)^{5} = 0.032$؛ وخطيًا $5/154 = 0.032$.
\textbf{4.} $6\times 0.32 = \qty{1.9}{\%}$.
\textbf{5.} $(1 - 1/154)^{5} = 0.968$: أي \qty{96.8}{\%} محفوظة.
\textbf{6.} عشرة أجيال عند $N_{e} = 53$: الزيادة $1 - (1 -
1/107)^{10} = 0.090$، والمجموع $F = 1 - 0.968\times 0.910 = 0.12$؛ وفقد
الملاءمة $6\times 1.2 = \qty{7}{\%}$؛ وذلك قرن.
\textbf{7.} أزواج الأصليين بالأصليين $(60/70)^{2} = 0.73$ من التزاوجات
ويحفظ خلفها $F \approx 0.03$ (وزيادة صغيرة)؛ وأما البقية
فعندها $F = 0$: فيكون متوسط $F \approx 0.73\times 0.035 = 0.026$، أي هبوط
قدره الربع في جيل واحد.
\textbf{8.} $F$ هوية بالنسب داخل الفرد، وأي خلف من تزاوج
بعيد يعيده إلى الصفر؛ وأما \omterm{thm:b3:conservation-biology:ne}{التغايرية الزيجية} فخاصة من خصائص
تواترات الأليلات، ولا تنتشر أليلات الوافدين إلا إذ
تُورَّث. وفي جمهرة لا تزال صغيرة، سيزيلها الانجراف ثانيةً
في غضون عشرات الأجيال، فلا بد من تكرار الإنقاذ أو من تكبير
الجمهرة.
\textbf{9.} $N_{e} = N$ عند تساوي الجنسين: أي 500 بالغ. ومن 70 عند $r =
0.05$: $\ln(500/70)/0.05 = 39$ سنة.
\textbf{10.} الجمهرة المثالية أحجام أسرها بواسونية، تباينها
يساوي وسطيها وهو اثنان؛ و$N_{e} = 4N/(2 + V_{k})$، فجعل كل زوج
يسهم بالقدر نفسه ($V_{k} = 0$) يعطي $N_{e} = 2N$: إذ لا تُفقد أي
سلالة بحظ العش.
\textbf{11.} اجمع نطافه واحفظها للتلقيح الاصطناعي،
واستعملها على عدة إناث غير قريبة عبر السنين --- والكلفة
أن أليلاته قد تهيمن على مجمع صغير، فلا بد من تسقيف حصته
قرب $1/N_{e}$؛ أو احفظ نسيجًا بالتبريد لاستنساخ لاحق أو
لإنتاج أعراس، وكلفة ذلك التأخير والخطر التقني. وألّا يُفعل
شيء يفقد أليلات السلالة إلى الأبد.
\textbf{12.} \qty{160}{} مليون دولار لنحو 430 طائرًا: أي نحو
\num{370000} للطائر. وهو غال لكل ببغاء، رخيص في مقابل نوع
--- وأقل من بضعة كيلومترات من طريق سريع، وهي المقارنة التي تجريها
الميزانيات فعلًا.
\textbf{13.} $(300/2000)^{0.25} = 0.62$: أي نحو 112 نوعًا من 180 محفوظة،
ودين انقراض قدره 68.
\textbf{14.} كل شظية $(100/2000)^{0.25}\times 180 = 85$ نوعًا.
ومختلفة تمامًا: حتى 255، وهو مستحيل فوق مجمع من 180؛
ومتطابقة: 85. والذي يقرر هو مقدار اختلاف موائل الشظايا
وكم نوعًا يقدر على الدوام في \qty{100}{km^2}؛ وواقعيًا
بين 85 و112.
\textbf{15.} $(320/2000)^{0.25}\times 180 = 114$، في مقابل 85 لثلاث
شظايا متطابقة معزولة: فالممر يدع إعادة الاستعمار
وتدفق المورّثات يجعلان ثلاث جمهرات صغيرة تسلك سلوك واحدة أكبر.
\textbf{16.} $N_{e} = 50$ عند تساوي الجنسين هو 25 زوجًا: أي \qty{1250}{km^2}.
و\qty{300}{km^2} كلها تحمل ستة أزواج، $N_{e} \approx 12$؛ والشظية
زوجين. ولا تكفي واحدة منهما: فالمفترس يُفقد مهما كان عدد
الأنواع، وحماية المساحة التي يحتاج إليها تحمي كل ما هو
أصغر --- أي إنه نوع مظلة.
\textbf{17.} $\ln(200/10)/0.03 = 100$ سنة: فالدين يُسدَّد على مدى
قرن، بعد قطع الغابة بزمن طويل.
\textbf{18.} الحواف تنزع الموئل الداخلي من كل شظية، فتكون
المساحة القابلة للاستعمال دون المرسومة ويكون الفقد أكبر مما
يتنبأ به القانون؛ وأما نسيج من نموّ ثانوي أو أسيجة شجرية تقدر بعض
الأنواع على استعماله فيضيف مساحة فعالة ويجعل الفقد أصغر.
\textbf{19.} $r = \ln(100/31)/8 = \qty{0.146}{yr^{-1}}$؛ وزمن المضاعفة
$\ln 2/0.146 = 4.7$ سنة.
\textbf{20.} $N(8) = 120/[1 + (120/31 - 1)\mathrm{e}^{-1.17}] = 120/(1 +
2.87\times 0.31) = 63$. والمرصود 100 هو بعينه قيمة النموذج
الأسّي --- إذ حُسب معدل النمو من هاتين النقطتين نفسيهما --- في
مقابل 63 في اللوجستي: ففي السنين الأولى لم تلقَ الذئاب
اعتمادًا على الكثافة --- أقاليم فارغة وأيائل وفيرة.
\textbf{21.} $\ln(4000/\num{17000})/20 = -0.072$: أي \qty{7}{\%} في السنة.
ومئة ذئب تأخذ \num{2000} أيل في السنة من قطيع متوسطه
\num{8000} هي الربع في السنة --- أي أكثر من كافية لسوق التراجع
وحدها، وإن كان الصيد خارج المتنزه والجفاف وافتراس الدببة
الرمادية للعجول قد شاركت فيه.
\textbf{22.} ستون فردًا ناضجًا، دون 250: مهدَّد بالانقراض تهديدًا حرجًا.
والمفترس عند $N_{e} \approx 12$ (بضع عشرات من الحيوانات): مهدَّد بالانقراض
تهديدًا حرجًا كذلك.
\textbf{23.} من $3\times 10^{9}$ إلى $1$ في 44 سنة: $r = -\ln(3\times
10^{9})/44 = -0.50$ في السنة، أي تنصّف كل سبعة عشر شهرًا. ولا
حصاد يزيل نصف المليارات كل سنة؛ بل كان الهبوط أبطأ في
أوله وأسرع كثيرًا في آخره، لأن متكاثرًا مستعمِرًا دون
كثافة عتبية يكف عن التكاثر بينما لم تعد مفترساته تُشبَع
--- فقد أحال \omterm{def:b3:conservation-biology:small}{أثر ألي} تراجعًا انهيارًا.
\textbf{24.} كانت الجمهرة عام 1850 بالمليارات، ومعدل نموها المقيس
موجبًا وتباينه صغيرًا، وتحليل القابلية للبقاء يسقط الحركيات
التي تُعطى له: فكان سيضع الخطر عند الصفر.
وكان سيفوته السوق الخارجي --- السكك الحديدية والتلغراف
اللذان أتاحا للصيادين تعقّب كل مبيت وإفراغه، وقطع الغابات --- وأن
هناك عتبة ألي في متكاثر مستعمِر، وليس واحد منهما في
ديموغرافيا نوع وافر؛ فالتحليل استقراء، وأسباب
الانقراض جديدة عادةً.
\textbf{25.} $N_{e} \approx 53$ اليوم؛ و$F \approx 0.12$ بعد عشرة أجيال
أخرى عند 60 (أي \qty{7}{\%} من الملاءمة)؛ والغابة المتقلّصة تبقي
نحو \qty{62}{\%} من أنواعها.
\end{solution}
